\chapter{变量定义与补充说明}
\label{chap:appa-variable-notes}

\section{核心变量定义}
\label{sec:appa-variable-definitions}

本附录汇总正文经验分析中频繁出现的核心变量，并补充其构造逻辑，便于读者与第五章的基准回归结果对照阅读。与稳健性结果有关的补充表格放在附录~\ref{chap:appb-robustness-tables}。表~\ref{tab:appa-core-variables} 展示了主要变量的经济含义、计量口径和预期方向；式~\eqref{eq:appa-monitoring-index} 则给出一个示例性的监督强度指数，用于说明变量标准化后的组合方式。

\begin{table}[htbp]
    \centering
    \caption{核心变量定义示例}
    \label{tab:appa-core-variables}
    \begin{tabular}{p{3.1cm}p{4.2cm}p{5.6cm}}
        \toprule
        变量名 & 计量方式 & 说明 \\
        \midrule
        shirking\_score & 1--5 的问卷指数 & 数值越大表示员工在工作时间从事非任务活动的频率越高。 \\
        monitor\_intensity & 标准化监督次数 & 根据月度抽查、系统留痕和主管回访加总后标准化。 \\
        incentive\_pay & 绩效工资占比 & 表示浮动收入在总薪酬中的比例，用于衡量激励强度。 \\
        task\_complexity & 工作复杂度评分 & 由任务多样性、知识密集度与跨部门协作要求共同构成。 \\
        \bottomrule
    \end{tabular}
\end{table}

为了展示附录中的公式、编号和交叉引用是否工作正常，下面给出一个简单的综合指数定义：
\begin{equation}
    \label{eq:appa-monitoring-index}
    M_i = \frac{1}{3}\left(z(\text{spotcheck}_i) + z(\text{logtrace}_i) + z(\text{managerreview}_i)\right).
\end{equation}

其中，$z(\cdot)$ 表示对原始指标进行标准化处理。正文若需要解释监督变量的构造过程，可直接引用附录的式~\eqref{eq:appa-monitoring-index} 和表~\ref{tab:appa-core-variables}。
